% ============================================================
%  sections/litreview.tex  –  Literature Review
% ============================================================

The international evidence consistently indicates a negative price elasticity of demand for cannabis, but estimates span a wide range depending on the behavioral margin examined, the measurement strategy, and the institutional setting in which transactions occur.  This section organizes the relevant literature along these dimensions (meta-analytic syntheses, micro-level measurement, crowdsourced transaction data, cross-country evidence with quality controls, the participation margin, legalization-era identification, and Mexico-specific analysis) and concludes by motivating the contribution of the present paper.

The broadest synthesis is the meta-analysis of illicit drug demand by \citet{gallet2014monkey}, which compiles hundreds of elasticity estimates across substances and reports a median price elasticity of demand for marijuana of $-0.33$.  Meta-regressions in that study reveal systematic heterogeneity: long-run elasticities tend to exceed short-run elasticities in absolute value, participation-margin estimates are generally more negative than intensive-margin estimates, and studies relying on hypothetical price scenarios produce substantially larger responses than those using observed transactions.  These patterns underscore the value of transaction-level data and motivate caution when extrapolating across settings.

Early micro evidence illustrates how the choice of data source shapes estimated responsiveness even within a single population.  Using survey data on students at the University of California, Los Angeles, \citet{nisbet1972ucla} estimate elasticities from both a hypothetical market-survey demand schedule and realized purchase records.  The double-log specification yields an elasticity of $-1.013$ for realized purchases versus $-0.365$ for the market survey; evaluated at the then-prevailing price of \$10 in a linear model, the analogous figures are $-1.51$ and $-0.51$.  The gap is directly relevant for illicit-market analysis: stated demand schedules may fail to capture search costs, legal risk, and other frictions, whereas realized transactions incorporate the constraints buyers actually face.  This early finding foreshadows a persistent methodological lesson in the literature, namely that the data-generating process matters as much as the econometric specification for recovering behaviorally meaningful demand parameters.

More recent work exploits transaction-level information from crowdsourced platforms and typically finds inelastic but economically meaningful responses.  \citet{davis2016crowd} use crowdsourced U.S.\ transaction data from \texttt{PriceOfWeed.com}, an online price-reporting platform, and estimate an OLS elasticity of $-0.689$.  To address potential price endogeneity, they employ instrumental-variables strategies using plausibly supply-driven shifters (electricity prices and geographic distance to major growing regions) and report IV estimates ranging from $-0.597$ to $-0.789$ depending on the instrument set.  The identification logic is that these cost shifters affect the retail price through the supply side but, conditional on controls, do not independently influence the quantity a given buyer demands.  Similarly, \citet{halcoussis2017geo} use \texttt{PriceOfWeed.com} data to model price and quantity jointly using geographic cost variation tied to distance from production areas, reporting a price elasticity of $-0.418$.  Extending the crowdsourced approach to a different national context, \citet{giommoni2018uk} analyze the British cannabis market using \texttt{PriceOfWeed.com} reports and document substantial regional price variation, providing further evidence that crowdsourced transaction data can generate economically informative estimates of illicit-market structure.  Together, these studies suggest that transaction-based elasticities for cannabis typically fall in the range of roughly $-0.4$ to $-0.8$, while highlighting that credibility depends on the defensibility of exclusion restrictions.

Survey-based transaction evidence from outside the United States points to comparable magnitudes and underscores the importance of controlling for product quality.  Using an online survey of South African cannabis users, \citet{riley2020southafrica} estimate a baseline conditional elasticity of $-0.501$.  Across alternative specifications that vary the controls for product differentiation, the coefficient ranges from $-0.493$ to $-0.613$.  This sensitivity to quality-related controls implies that, in markets where product heterogeneity is pronounced, omitting quality proxies risks biasing the estimated demand response.  The concern is directly relevant here because, as shown below, seedless and seeded marijuana command sharply different prices in the Mexican survey data.  \citet{benlakhdar2016potency} reinforce this point using data from over 250 regular cannabis users in France, estimating short-run price elasticities between $-1.7$ and $-2.1$.  Notably, controlling for real or perceived potency has little effect on the magnitude of the price response in their sample, suggesting that while quality matters for willingness to pay, the price elasticity itself may be relatively robust to potency measurement once product heterogeneity is accounted for.

A related strand of the literature estimates elasticities on the extensive (participation) margin.  \citet{pacula2000youth} show that youth participation elasticities are sensitive to trend specification and the inclusion of mediators, reinforcing the distinction between participation and intensive-margin responses.  \citet{amlung2019substitution} complement this work by documenting that demand for illegal cannabis is relatively more elastic than demand for legal cannabis, implying that consumers are willing to switch sources as relative prices shift.

Recent legalization in several U.S.\ states and in Canada has created new opportunities for identification via fiscal-policy and regulatory variation.  \citet{han2025commercialization} study adult past-30-day use in states with recreational commercialization and exploit tax rates as an instrument for retail prices, reporting implied elasticities between $-0.66$ and $-0.59$.  \citet{dodd2025canada} provide the first post-legalization elasticity estimate for Canada's regulated market, using sales data from British Columbia and reporting a price elasticity of approximately $-1.4$, substantially more elastic than most estimates from illicit settings.  Their approach illustrates the appeal of regulated-market data, where prices, quantities, and product attributes are reliably observed, but the resulting estimates may not transfer directly to illicit markets where administered tax variation is absent, product quality is heterogeneous, and buyers face additional legal and search-cost frictions.

Mexico-specific work has engaged with the taxation question but largely without transaction-level evidence capable of identifying a domestic price elasticity.  \citet{perezromero2014elementos} develop a microeconomic framework for analyzing cannabis legalization and indirect taxation.  A central insight of their analysis is that the relevant price faced by consumers in an illicit market should be understood as a ``full price'' encompassing not only the monetary cost but also legal risk, search costs, and other transaction frictions associated with purchasing from an illegal supplier.  They further discuss how substitution and complementarity with legal substances, notably alcohol and tobacco, may shape the effects of legalization on consumption patterns.  At the same time, their discussion highlights a binding constraint: in the absence of reliable price and consumption data from the Mexican market, direct estimation of demand parameters has not been feasible.

Two points follow from this literature.  First, cannabis demand slopes downward and is often less than unit-elastic, but estimates vary across behavioral margins, measurement strategies, and institutional settings \citep{gallet2014monkey}.  Second, credible estimation requires explicit attention to price endogeneity and product differentiation \citep{davis2016crowd,halcoussis2017geo,riley2020southafrica,benlakhdar2016potency}.  No study to date has estimated the transaction-level price elasticity of marijuana demand using purchase records from Mexico.  This paper fills that gap with a raw crowdsourced survey of 4{,}775 purchase records that jointly observe price, quantity, perceived quality, and fine-grained geography \citep{bejaranoromero2020dataset}; after cleaning, the analytic sample retains 3{,}293 transactions.  The log--log demand framework controls for quality tiers, captures the sharp price wedge between seedless and seeded product documented in the data, and absorbs unobserved geography and timing through fixed effects, with the preferred specification using state-by-month fixed effects.

Price endogeneity is a first-order concern in any demand analysis, and illicit markets pose additional identification challenges because standard supply-side instruments (taxes, regulated input costs) are unavailable \citep{han2025commercialization,dodd2025canada}.  We address this concern along two complementary dimensions.  First, our baseline specifications systematically vary controls, fixed-effects structures, and sample restrictions, allowing us to assess coefficient stability and the sensitivity of the estimated elasticity to potential confounders; we formalize this approach via the \citet{oster2019unobservable} bounding procedure.  Second, we implement instrumental-variables strategies motivated by the geography-based cost shifters used by \citet{davis2016crowd} and \citet{halcoussis2017geo}, constructing instruments based on Open Source Routing Machine (OSRM) driving times over OSM road data to major production hubs, a one-month-lag OSM driving-time-weighted seizure index drawn from M\'{e}xico Unido contra la Delincuencia (MUCD) municipality-month enforcement data \citep{openstreetmap2026,luxen2011osrm,mucd2019seizures}, and a Hausman-type leave-one-out mean price.  We report these IV results primarily as diagnostics: the Hausman-type instrument provides a strong internal robustness check under its leave-one-out exclusion restriction, while the external geography-based instruments require weak-IV diagnostics because their first stages are limited.  This transparent, assumption-explicit approach provides an empirical foundation for the ongoing policy discussion on cannabis regulation and taxation in Mexico.
